\subsection{Operation-Specific Construction Rules}
\label{app:canonical-axis-rules}
\label{app:ordering-rules}
\label{app:attribute-construction-rules}
Let $I_X$ list input $X$'s logical axes in their original order, using the
LHS axis identities. The following rules define retained axes, ordering
checks, and attribute construction under D1--D4. Compute $P_v$ from inputs
to root, propagate $Req_v$ from root to inputs starting with
$Req_{\mathrm{root}}=F$, and construct attributes from inputs to root after
all checks succeed. A requirement must contain distinct axes from $P_v$;
otherwise $C_v$ is false. All child requirements are checked recursively.

\noindent\textbf{Input.}
For input $X$, set $P_X=\operatorname{set}(I_X)$ and
$C_X=[Req_X\preceq I_X]$. Its output order is
$\operatorname{Ord}_X=I_X$; no attributes are needed.

\noindent\textbf{Permutation.}
Set $P_v=P_u$, pass $Req_u=[]$, and take $C_v=\mathrm{true}$.
Construct $\operatorname{Ord}_v$ by placing $Req_v$ first and appending
the remaining child axes in their current order. The permutation attribute
lists their positions in the child.

\noindent\textbf{Unary summation.}
Set $P_v=P_u\cap M_v$, pass $Req_u=Req_v$, and take
$C_v=\mathrm{true}$. The output order is
$\operatorname{Ord}_v=\operatorname{Ord}_u|_{P_v}$.
The reduction attribute lists the positions of $P_u\setminus P_v$
in the child.

\noindent\textbf{Contraction.}
For children $\ell,r$, define the batch, contraction, left-free, and
right-free axis sets by
\begin{equation}
\begin{aligned}
 B_v&=P_\ell\cap P_r\cap M_v,&
 K_v&=(P_\ell\cap P_r)\setminus M_v,\\
 LF_v&=P_\ell\setminus P_r,&
 RF_v&=P_r\setminus P_\ell.
\end{aligned}
\label{eq:canonical-roles}
\end{equation}
Set $P_v=B_v\cup LF_v\cup RF_v$ and pass
$Req_\ell=Req_v|_{LF_v}$ and $Req_r=Req_v|_{RF_v}$.
The local check $C_v$ is
\begin{equation}
 \left[Req_v=(Req_v|_{B_v})\Vert(Req_v|_{LF_v})\Vert(Req_v|_{RF_v})\right].
 \label{eq:block-order}
\end{equation}
Form the batch sequence $b$ by placing $Req_v|_{B_v}$ first, then appending
the remaining batch axes in a fixed order. List the contraction axes
$K_v$ in a fixed order $k$. Use the same sequences for both operands.
The four attribute lists are
\begin{equation}
\begin{aligned}
 b_\ell&=\operatorname{pos}(\operatorname{Ord}_\ell,b),&
 b_r&=\operatorname{pos}(\operatorname{Ord}_r,b),\\
 k_\ell&=\operatorname{pos}(\operatorname{Ord}_\ell,k),&
 k_r&=\operatorname{pos}(\operatorname{Ord}_r,k).
\end{aligned}
\label{eq:formal-attributes}
\end{equation}
The output lists batch axes first, followed by free axes in their operand
order:
\begin{equation}
 \operatorname{Ord}_v=b\Vert
 (\operatorname{Ord}_\ell|_{LF_v})\Vert
 (\operatorname{Ord}_r|_{RF_v}).
 \label{eq:constructed-order}
\end{equation}
Only axes whose supporting inputs are all inside the subtree and which
are absent from the final output may be removed. Contraction removes
such shared axes; unary summation can remove any such axis in its input.

\subsection{Mandatory and Retained Axes: A Worked Example}
\label{app:canonical-axes-example}
Consider a rewrite that moves a row sum from an input to the result of
batched matrix multiplication:
\begin{equation}
 \operatorname{dot}(\operatorname{sum}_i(A),B)
 \;\Rightarrow\;
 \operatorname{sum}_i(\operatorname{dot}(A,B)).
\end{equation}
Here $\operatorname{dot}$ preserves the shared batch axis $b$ and contracts
axis $k$, while $\operatorname{sum}_i$ removes row axis $i$. The LHS sums
the rows of $A$ before multiplication; the RHS multiplies first and sums
the resulting rows. We apply Section~\ref{sec:canonical-axes} to the RHS,
where multiplication must keep $i$ until the final summation even though
it is not needed in the final output.

\subsubsection{Inputs and Logical Axes}
Let $A,B$ have shapes $(s,m,n)$ and $(s,n,p)$, with all extents at least two.
Their ordered logical axes are
\begin{equation}
 I_A=(b,i,k),\qquad I_B=(b,k,j).
\end{equation}
Axes $b,k$ connect the inputs, while $i,j$ occur in only one input each:
\begin{equation}
 \sigma(b)=\sigma(k)=\{A,B\},\qquad
 \sigma(i)=\{A\},\qquad \sigma(j)=\{B\}.
\end{equation}
In indexed form, the source $L$, RHS intermediate $T$, and RHS output $Y$ are
\begin{equation}
\begin{aligned}
 L_{bj}&=\sum_k\left(\sum_i A_{bik}\right)B_{bkj},\\
 T_{bij}&=\sum_k A_{bik}B_{bkj},\qquad
 Y_{bj}=\sum_i T_{bij}.
\end{aligned}
\end{equation}
Both computations have output order $F=(b,j)$ and agree under exact real
arithmetic. The RHS nodes are $A,B,T,Y$.

\subsubsection{Mandatory Axes and Axes Kept}
At input $A$, axis $b$ is required in the output and $k$ must connect to $B$.
In particular, $\sigma(k)=\{A,B\}\nsubseteq U_A=\{A\}$: an input
containing $k$ lies outside this subtree. Axis $i$ is not mandatory, but
an input node keeps all its axes. Thus $M_A=\{b,k\}$ while
$P_A=\{b,i,k\}$.

At $T$, both inputs are inside the subtree, so only the final output axes
are mandatory: $M_T=\{b,j\}$. Multiplication sums out $k$, but cannot
sum out $i$, which appears only in its left operand. It must keep $i$ as
a free axis, giving
\begin{equation}
 M_T=\{b,j\}\subsetneq P_T=\{b,i,j\}.
\end{equation}
The final unary summation removes $i$. Table~\ref{tab:matmul-axis-sets}
lists the sets at every RHS node.
\begin{table}[!htbp]
\caption{Mandatory and retained axes for batched multiplication followed
by a row sum.}
\label{tab:matmul-axis-sets}
\begin{tabularx}{\linewidth}{@{}lYll@{}}\toprule
Node & Inputs in subtree & Mandatory axes & Axes kept\\\midrule
$A$ & $\{A\}$ & $\{b,k\}$ & $\{b,i,k\}$\\
$B$ & $\{B\}$ & $\{b,k,j\}$ & $\{b,k,j\}$\\
$T=AB$ & $\{A,B\}$ & $\{b,j\}$ & $\{b,i,j\}$\\
$Y=\operatorname{sum}_i(T)$ & $\{A,B\}$ & $\{b,j\}$ & $\{b,j\}$\\\bottomrule
\end{tabularx}
\end{table}

\subsubsection{Batch, Contraction, and Free Axes}
At $T$, the operands share $b,k$. Axis $b$ is mandatory and stays as batch;
$k$ is not mandatory and is summed out. Axes $i,j$ remain free.
Table~\ref{tab:matmul-axis-roles} lists the four roles.
\begin{table}[!htbp]
\caption{Axis roles at the multiplication node $T$.}
\label{tab:matmul-axis-roles}
\begin{tabularx}{\linewidth}{@{}llY@{}}\toprule
Role & Axis set & Reason\\\midrule
Batch & $B_T=\{b\}$ & Shared and required in the final output.\\
Contraction & $K_T=\{k\}$ & Shared and no longer needed.\\
Left-free & $LF_T=\{i\}$ & Present only in $A$; multiplication must keep it.\\
Right-free & $RF_T=\{j\}$ & Present only in $B$; required in the final output.\\\bottomrule
\end{tabularx}
\end{table}
Thus $P_T=B_T\cup LF_T\cup RF_T=\{b,i,j\}$. At the root, unary
summation keeps $P_Y=P_T\cap M_Y=\{b,j\}=\operatorname{set}(F)$.
The extra axis $i$ survives multiplication only because that operation
cannot remove it.
